\chapter{Variables and expressions: teaching the model to calculate}

The next step in parametric modelling is to write design intent into \textbf{variables}: gather length,
width, hole diameter and pattern count into a single spreadsheet, and let every dimension reference it.
Change one place and the whole design updates. That is what the ``Variables'' tab (the Spreadsheet
workbench) is for.

\section{Opening and creating}

\begin{enumerate}
  \item Switch to the \ui{Variables} tab at the bottom; the empty state shows \ui{No spreadsheet yet} and
        \ui{Create spreadsheet};
  \item Tap \ui{Create spreadsheet} (\cmd{Spreadsheet\_CreateSheet});
  \item The sheet replaces the 3D view; a \ui{Spreadsheet} object appears in the tree (default name
        \file{Spreadsheet}).
\end{enumerate}

\section{Cells and input}

\begin{itemize}
  \item \textbf{Select}: tap a cell. \textbf{Range}: touch and hold to anchor the start, then tap the
        opposite corner; you get an area such as \file{A1:C4};
  \item \textbf{Edit}: double-tap a cell, or use the \ui{Edit} button on the right of the formula bar
        above the grid;
  \item There are three kinds of input:
    \begin{itemize}
      \item Text: just type it;
      \item A number with a unit: \file{25 mm}, \file{30 °};
      \item A formula: starts with \file{=}, for example \file{=A2*2} or \file{=Sheet2.B3};
    \end{itemize}
  \item The left side of the formula bar shows the current selection and the right side shows its content
        (with the cell's alias in parentheses if it has one, for example ``B3 (Thickness)'').
\end{itemize}

\section{Aliases: naming parameters}

\begin{enumerate}
  \item Select a cell;
  \item Run \ui{Set alias} (\cmd{Spreadsheet\_SetAlias});
  \item Type a name (for example \file{Thickness}).
\end{enumerate}

After that you can write \file{=Sheet.Thickness} in \textbf{any} numeric property (\file{Sheet} is the
spreadsheet's object name and can be changed in the tree). If the alias is unique in the whole document,
you can also shorten it to \file{=Thickness}.

\section{Expressions: usable anywhere}

Property rows (including values in the feature dialog, sketch dimensions, FEM loads, assembly offsets,
and so on) all accept expressions:

\begin{itemize}
  \item Enter a plain ``number + unit'': it is stored as a value;
  \item Enter anything else (an operation or a reference): it is stored as an \textbf{expression}, and the
        property row shows a blue \file{=...};
  \item Lengths and angles are both supported; vector components (for example the X/Y/Z of a placement)
        follow the same rules.
\end{itemize}

\begin{lstlisting}
=Sheet.Thickness            // refer to a variable
=Sheet.Width - 2 * 3 mm     // leave a margin
=10 mm + 5 mm               // arithmetic directly
\end{lstlisting}

\section{Formatting and data exchange}

The toolbar also provides six alignments, bold/italic/underline, merge/split cells, column width and row
height, insert/delete rows and columns, and \ui{Import}/\ui{Export} of CSV/TSV (the separator follows the
extension: comma, semicolon or tab).

\section{Unit system}

\menupath{Edit → Preferences → Unit system} switches FreeCAD's unit scheme and number of decimal places.
It affects \textbf{display and input parsing}; internally the app always uses millimetres and kilograms,
so switching leaves existing values unchanged---only the way they are displayed changes.

\begin{tip}[Treat the parameter sheet as a design specification]
Build a sheet for every model: \file{Length}, \file{Width}, \file{Height}, \file{HoleDia},
\file{PatternCount}, \ldots{} and write comment rows beside them explaining what each one means. When you
come back to change the design three months later, you will thank yourself.
\end{tip}

\begin{practice}[Exercise: make the bracket configurable]
\begin{enumerate}
  \item On the Variables tab, create a sheet; put \file{Width} in A1 (value 60 mm), \file{Height} in A2
        (40 mm) and \file{Hole} in A3 (10 mm);
  \item Select B1/B2/B3 and give them the aliases \file{Width}, \file{Height} and \file{Hole};
  \item Go back to the bracket sketch and change the two dimensions to \file{=Sheet.Width} and
        \file{=Sheet.Height};
  \item Change the hole diameter to \file{=Sheet.Hole};
  \item Return to the Variables tab and change B1 to 100 mm---the bracket widens automatically, and the
        holes and fillets update with it.
\end{enumerate}
\end{practice}

\section{Limitations}

\begin{itemize}
  \item Formula entry is a pop-up plus the numeric keypad; there is no function-completion drop-down as in
        the desktop version;
  \item The cell type (text/number/formula) is decided from what you type; there is no explicit type
        setting;
  \item Large sheets have no filtering, sorting or charts; this is essentially a ``parameter table'', not
        Excel.
\end{itemize}
